\chapter{TỔNG QUAN ĐỀ TÀI}\label{Chapter1}

\section*{Tóm tắt chương}

Chương này trình bày bối cảnh, động lực nghiên cứu, các công trình liên quan, mục tiêu, đối tượng, phạm vi, phương pháp nghiên cứu và cấu trúc khóa luận. Trọng tâm của chương là lý giải vì sao bài toán IDS hiện nay không chỉ là bài toán huấn luyện một mô hình chính xác, mà còn là bài toán học trong môi trường phân tán, dữ liệu Non-IID, mô hình không đồng nhất và lớp tấn công thay đổi theo thời gian.

\section{Lý do chọn đề tài}

Trong môi trường số hiện nay, hệ thống mạng đã trải rộng từ trung tâm dữ liệu sang điện toán đám mây, điện toán biên, thiết bị IoT, hệ thống công nghiệp kết nối, camera thông minh và nhiều loại thiết bị đầu cuối khác nhau. Khi hạ tầng mở rộng theo hướng đó, mỗi thành phần không chỉ sinh dữ liệu và tạo thêm kết nối, mà còn đồng thời mở rộng bề mặt tấn công của toàn bộ hệ thống. Các báo cáo an ninh mạng gần đây cho thấy lưu lượng độc hại nhắm vào IoT tiếp tục tăng mạnh, botnet quy mô lớn vẫn hình thành từ các thiết bị cấu hình yếu, còn tấn công từ chối dịch vụ phân tán (Distributed Denial of Service attacks - DDoS attacks) tiếp tục khai thác chính những thiết bị đặt ở rìa mạng \cite{sonicwall2024report,netgear2024report,hackernews2026report}. Vì vậy, hệ thống giám sát không còn chỉ cần phát hiện một vài mẫu bất thường đơn lẻ, mà phải theo dõi một không gian hoạt động rộng hơn, không đồng nhất hơn và biến động liên tục theo thời gian.
\bigbreak\noindent
Trong bối cảnh đó, hệ thống phát hiện xâm nhập (Intrusion Detection System - IDS) giữ vai trò như một lớp quan sát trung gian giữa hạ tầng mạng và các biện pháp bảo mật phía sau. IDS hiệu quả không chỉ tạo ra cảnh báo, mà còn hỗ trợ người vận hành nhận biết cuộc tấn công bắt đầu từ đâu, lưu lượng nào đang lệch khỏi hành vi thông thường, mẫu bất thường đó thuộc lớp đã biết hay là biến thể mới, và quyết định của mô hình có đủ cơ sở để tin cậy hay không. Các phương pháp IDS dựa trên chữ ký (signature-based) vẫn hữu ích với những mẫu tấn công đã được mô tả rõ, nhưng lợi thế này giảm dần khi biến thể tấn công thay đổi nhanh hơn chu kỳ cập nhật luật. Chính từ giới hạn đó, học máy và học sâu (Deep Learning - DL) được đưa vào IDS để học trực tiếp từ dữ liệu lưu lượng, nơi mối quan hệ giữa cổng, giao thức, số gói, số byte, thời lượng phiên và các đặc trưng tổng hợp thường phức tạp hơn nhiều so với những gì luật thủ công có thể bao quát ổn định.
\bigbreak\noindent
Khó khăn lớn nằm ở chỗ dữ liệu an ninh mạng không dễ được tập trung về một nơi để huấn luyện. Dữ liệu mạng nội bộ thường chứa thông tin nhạy cảm về hạ tầng, hành vi người dùng, cấu trúc phân vùng hệ thống và tín hiệu vận hành. Ở các môi trường công nghiệp, năng lượng, y tế hoặc tài chính, yêu cầu này còn chặt hơn do các ràng buộc vận hành và pháp lý như GDPR \cite{gdpr}. Ngay cả khi tạm gác yếu tố quyền riêng tư, việc truyền dữ liệu thô khối lượng lớn từ nhiều điểm biên về một trung tâm vẫn làm tăng chi phí lưu trữ, băng thông và thời gian xử lý. Vì thế, bài toán IDS rơi vào một mâu thuẫn thực tế: hệ thống muốn học tốt thì cần dữ liệu phong phú từ nhiều nguồn, nhưng những dữ liệu có giá trị nhất lại thường là những dữ liệu khó chia sẻ nhất.
\bigbreak\noindent
Học liên kết (Federated Learning - FL) xuất hiện như một hướng giải quyết trực tiếp cho mâu thuẫn đó, khi mỗi client giữ dữ liệu tại chỗ, tự huấn luyện trên dữ liệu cục bộ và chỉ chia sẻ tri thức ở mức mô hình \cite{fedavg}. Các nghiên cứu FL-IDS, DeepFed, FELIDS và nhiều hệ thống tương tự đã cho thấy khả năng kết hợp tri thức từ nhiều miền mạng khác nhau trong khi dữ liệu huấn luyện vẫn nằm ở từng client riêng lẻ \cite{li2020deepfed,FRIHA202217felids,bhavsar2024flids,popoola2021flids}. Tuy nhiên, FL truyền thống vẫn thường dựa trên trao đổi tham số hoặc gradient mô hình, khiến chi phí truyền thông tăng theo kích thước mô hình và số vòng học. Bên cạnh đó, cập nhật mô hình cũng có thể làm lộ thông tin huấn luyện thông qua các hình thức suy diễn hoặc tái tạo mẫu \cite{zhu2019dlg,zhao2020idlg,geiping2020inverting}. Trong môi trường IDS, nơi dữ liệu mạng phản ánh hoạt động nội bộ và hành vi người dùng, bề mặt rủi ro này cần được xem là một hạn chế quan trọng.
\bigbreak\noindent
Một giới hạn khác của FL dựa trên trọng số là áp lực đồng nhất kiến trúc giữa các client. Trong triển khai thực tế, mỗi miền mạng có thể có tài nguyên tính toán, yêu cầu độ trễ và đặc điểm dữ liệu khác nhau; có nơi phù hợp với mô hình nhỏ, có nơi dùng GRU để khai thác đặc trưng tuần tự, có nơi lại cần kiến trúc sâu hơn để tăng năng lực biểu diễn. Nếu tất cả client phải dùng chung một kiến trúc, hệ thống buộc phải chọn một phương án dung hòa, trong khi phương án đó không nhất thiết tối ưu cho từng client riêng lẻ. Học chắt lọc liên kết (Federated Distillation - FD) phù hợp hơn với tình huống này, vì client có thể truyền logits, soft label, prototype hoặc biểu diễn trung gian thay vì truyền toàn bộ trọng số mô hình \cite{hinton2015kd,jeong2018federateddistillation,li2019fedmd,li2024fdsurvey}. Khi các client chỉ cần gặp nhau ở không gian dự đoán, hệ thống có thể hỗ trợ mô hình không đồng nhất tốt hơn và giảm lượng thông tin phải truyền qua mạng.
\bigbreak\noindent
Mặt khác, nhiều nghiên cứu đã chỉ ra mối nguy tiềm ẩn trong việc trao đổi trọng số. Từ các bản cập nhật trọng số của các client, server hoặc một bên tấn công man-in-the-middle có thể thực hiện suy ngược mô hình và xấp xỉ được dữ liệu đã được dùng để huấn luyện mô hình đó. Đây là kiểu tấn công model inversion đã được nhiều nghiên cứu chỉ ra: Deep Leakage from Gradients (DLG) \cite{zhu2019dlg} và improved Deep Leakage from Gradients (iDLG) \cite{zhao2020idlg} có thể suy ngược các mẫu dữ liệu huấn luyện của một mô hình chỉ từ gradients của mô hình đó, trong khi \cite{geiping2020inverting} trực tiếp minh chứng cho kiểu tấn công này trong môi trường học máy liên kết, cũng như nhiều nghiên cứu khác \cite{melis2019mia, pyrgelis2017knockknock, ateniese2015inference, ganju2018inference} lợi dụng tính liên kết của FL để khai thác dữ liệu riêng tư của các client tham gia.
\bigbreak\noindent
Bên cạnh đó, nhiều hệ thống học máy liên kết vẫn phụ thuộc vào một server trung tâm. Chưa kể tới mối nguy của inference attack khi server có thể trực tiếp truy cập từng bản cập nhật của client, trong môi trường tập trung như vậy, vận hành của quá trình học hoàn toàn phụ thuộc vào sự ổn định của server. Từ đó, các nghiên cứu học máy liên kết phân tán như BALANCE \cite{fang2024balance} hay DFL-Dual \cite{sun2024dfldual} ra đời. Trong môi trường phi tập trung, client có thể luân phiên điều phối hoặc phối hợp ngang hàng, qua đó giảm sự phụ thuộc tuyệt đối vào server trung tâm \cite{braintorrent,beltran2023dflsurvey}. Nhiều nghiên cứu học máy liên kết cho IDS cũng đã ứng dụng mô hình phân tán để khắc phục nhược điểm này \cite{nakip2023dfl-ids, friha20232dfl-ids}.
\bigbreak\noindent
Từ các phân tích trên, khóa luận hướng đến một hệ thống IDS học cộng tác phù hợp hơn với môi trường triển khai thật, nơi dữ liệu phân tán, phân phối Non-IID, thiết bị không đồng nhất, chi phí truyền thông cần được kiểm soát và tri thức chia sẻ có thể sai lệch. Hệ thống này kết hợp FD phi tập trung (Decentralized FD - DFD), hỗ trợ mô hình không đồng nhất, đánh giá độ tin cậy bằng energy, biểu diễn bất định bằng EKD và học tiệm tiến (Class-Incremental Learning - CIL) để xử lý sự xuất hiện của lớp tấn công mới theo thời gian. Do đó, trọng tâm của đề tài không chỉ là xây dựng một mô hình IDS chính xác trong điều kiện lý tưởng, mà là thiết kế một cơ chế học liên tục và cộng tác có khả năng vận hành ổn định hơn trong bối cảnh dữ liệu an ninh mạng thực tế.

\section{Các nghiên cứu liên quan}

\subsection{Hệ thống phát hiện xâm nhập dựa trên học sâu}

Các hệ thống IDS dựa trên DL ra đời từ nhu cầu vượt qua giới hạn của hệ thống dựa trên chữ ký và của các mô hình học máy cổ điển vốn phụ thuộc nhiều vào đặc trưng thủ công. Thay vì phải thiết kế luật riêng cho từng dạng tấn công, mô hình DL học trực tiếp từ dữ liệu lưu lượng và tự xây dựng biểu diễn để phân biệt hành vi bình thường với hành vi bất thường. Cách tiếp cận này càng trở nên quan trọng khi lưu lượng mạng hiện đại có cấu trúc phức tạp hơn nhiều so với các bộ dữ liệu an ninh đời đầu, bởi một phiên mạng giờ đây không chỉ mang thông tin về cổng hay giao thức, mà còn phản ánh hành vi theo thời gian, quan hệ giữa chiều vào và chiều ra, cường độ truyền tải, cùng các mẫu phối hợp giữa nhiều đặc trưng cùng lúc.
\bigbreak\noindent
Trong nhóm mô hình DL, các mạng dense nhiều tầng thường được chọn khi dữ liệu đầu vào đã ở dạng đặc trưng tabular ổn định, nhờ ưu điểm dễ hiện thực, suy luận nhanh và thuận tiện khi làm mốc so sánh cơ sở. Khi bài toán cần khai thác phụ thuộc theo chuỗi hoặc theo thứ tự xuất hiện của tín hiệu, GRU và LSTM lại trở thành lựa chọn phù hợp hơn vì chúng mô hình hóa được quan hệ thời gian giữa các bước dữ liệu. Trong khi đó, với những hướng cần tận dụng khả năng trích xuất mẫu cục bộ hoặc biểu diễn kết hợp giữa không gian và thời gian, các kiến trúc lai như CNN--RNN hoặc DCNN--BiLSTM thường được sử dụng. Dù vậy, cần nhấn mạnh rằng không có một kiến trúc duy nhất tối ưu cho mọi miền mạng: một mô hình đạt kết quả tốt trên tập dữ liệu này chưa chắc giữ nguyên lợi thế trên tập dữ liệu khác, bởi phân phối đặc trưng, số lớp, mức mất cân bằng và cách trích xuất đặc trưng đầu vào đều có thể thay đổi theo từng bối cảnh dữ liệu.
\bigbreak\noindent
Đặt trong bối cảnh của khóa luận này, CICIoT2023 là một tập dữ liệu phù hợp để khảo sát vì phản ánh được môi trường IoT, có số lượng mẫu lớn và bao gồm nhiều lớp tấn công khác nhau. Nhờ đó, nó hữu ích cho việc đánh giá khi đặt mô hình trước nhiều tình huống hơn so với các bộ dữ liệu nhỏ và đơn giản. Thế nhưng, chính quy mô lớn ấy lại làm lộ ra một giới hạn quan trọng của IDS huấn luyện tập trung: dữ liệu càng lớn và càng đến từ nhiều nguồn thì càng khó gom về một nơi để duy trì một quy trình huấn luyện duy nhất. Vì thế, ngay cả khi DL cho thấy năng lực biểu diễn tốt, cách tổ chức huấn luyện vẫn là yếu tố quyết định khả năng ứng dụng trong hệ thống thật.
\bigbreak\noindent
Một thách thức khác của IDS dựa trên DL là mất cân bằng lớp. Trong lưu lượng mạng thực tế, dữ liệu bình thường thường chiếm tỷ lệ áp đảo, trong khi nhiều loại tấn công lại xuất hiện với tần suất rất thấp. Hệ quả là một mô hình có thể đạt độ chính xác cao nhờ học tốt lớp lớn nhưng vẫn bỏ sót các lớp tấn công hiếm, và đó cũng là lý do các chỉ số như precision, recall và F1-score trở nên quan trọng hơn nhiều so với việc chỉ nhìn vào accuracy đơn lẻ. Vấn đề còn phức tạp hơn khi dữ liệu được thu từ nhiều miền mạng khác nhau, bởi mỗi miền lại mang một kiểu mất cân bằng riêng: một client có thể quan sát nhiều mẫu quét cổng, client khác lại thường xuyên thấy lưu lượng DDoS, trong khi client còn lại chỉ có số ít mẫu thuộc một vài lớp hiếm. Chính sự chênh lệch này tạo nền cho những vấn đề của FL và FD được bàn ở các phần sau.
\bigbreak\noindent
Nhìn ở cấp độ hệ thống, các mô hình DL cho IDS cần đáp ứng đồng thời ba yêu cầu: phát hiện được các mẫu đã biết, thích nghi được với thay đổi của dữ liệu và vận hành được trên hạ tầng không đồng nhất. Xuất phát từ những yêu cầu đó, khóa luận không giới hạn ở một mô hình đơn lẻ, mà xem dense, GRU và các kiến trúc sâu hơn như những đại diện cho các lựa chọn triển khai khác nhau. Việc đặt nhiều kiến trúc vào cùng một bộ khung, nhờ vậy, cho phép đánh giá không chỉ hiệu năng phát hiện, mà còn cả khả năng phối hợp giữa các mô hình không đồng nhất.

\subsection{Học máy liên kết trong phát hiện xâm nhập}

Federated Learning là một bước phát triển gần như tự nhiên khi IDS chuyển từ môi trường dữ liệu tập trung sang môi trường dữ liệu phân tán. Trong nhóm này, FedAvg là nền tảng được sử dụng rộng rãi nhất \cite{fedavg}. Ý tưởng chính của FedAvg là mỗi client huấn luyện cục bộ trên dữ liệu riêng, sau đó gửi tham số mô hình về server để lấy trung bình có trọng số theo số lượng mẫu. Mô hình toàn cục được tạo ra từ đó sẽ được phát lại cho các client để tiếp tục vòng học kế tiếp. Nhờ quy trình rõ ràng và dễ triển khai, FedAvg thường được xem như điểm khởi đầu cho nhiều nghiên cứu FL-IDS.
\bigbreak\noindent
Đối với phát hiện xâm nhập, lợi ích của FL nằm ở chỗ mỗi đơn vị tham gia có thể giữ log mạng và đặc trưng lưu lượng trong miền cục bộ của mình. Điều này giúp giảm nhu cầu chuyển dữ liệu thô xuyên tổ chức, đồng thời giảm rủi ro lộ thông tin vận hành nhạy cảm. Các hướng như DeepFed, FELIDS và nhiều hệ thống FL-IDS khác đã cho thấy mô hình FL có thể cải thiện hiệu năng phát hiện trong môi trường công nghiệp, IoT và edge khi so với việc huấn luyện riêng lẻ trên từng client \cite{li2020deepfed,FRIHA202217felids,bhavsar2024flids,popoola2021flids}. Nói gọn hơn, giá trị cốt lõi của FL trong bối cảnh này là khả năng khai thác tri thức phân tán: mỗi client chỉ nhìn thấy một phần của môi trường mạng, nhưng mô hình toàn cục có thể học từ sự kết hợp của nhiều phần như vậy.
\bigbreak\noindent
Tuy nhiên, chính trong IDS thì các hạn chế của FedAvg cũng bộc lộ khá nhanh. Dữ liệu giữa các client thường Non-IID, nghĩa là không cùng phân phối, không cùng tỷ lệ lớp và thậm chí không cùng tập lớp quan sát nổi bật. Khi một client chủ yếu thấy một số hành vi tấn công nhất định, còn client khác lại chủ yếu thấy hành vi bình thường hoặc một nhóm tấn công khác, hướng tối ưu cục bộ của họ sẽ khác nhau rõ rệt. Nếu server vẫn lấy trung bình các tham số như thể mọi client đang học trên cùng một phân phối, mô hình toàn cục sẽ dễ dao động hoặc hội tụ chậm. Hiện tượng này thường được gọi là client drift trong FL trên dữ liệu không đồng nhất.
\bigbreak\noindent
Để giảm client drift, nhiều biến thể của FL đã được đề xuất. FedProx thêm một ràng buộc để mô hình cục bộ không trôi quá xa khỏi mô hình toàn cục \cite{li2020fedprox}, phù hợp khi muốn duy trì tính ổn định của quá trình học mà không thay đổi hoàn toàn quy trình FedAvg. FedDyn đi xa hơn bằng cách điều chỉnh mục tiêu tối ưu theo cách động để giảm lệch giữa tối ưu cục bộ và tối ưu toàn cục \cite{acar2021feddyn}. Trong khi đó, FedMLB khai thác tri thức đa mức giữa mô hình cục bộ và toàn cục để học ổn định hơn trên dữ liệu không đồng nhất \cite{kim2022fedmlb}. Nhìn chung, các hướng này cho thấy bài toán cốt lõi của FL-IDS không chỉ là giữ dữ liệu tại nguồn, mà còn là làm sao để tri thức từ nhiều client thực sự dung hòa được với nhau.
\bigbreak\noindent
Dù cải thiện được phần nào khả năng học trên dữ liệu Non-IID, FL trọng số vẫn giữ lại ba hạn chế nền tảng. Thứ nhất là chi phí truyền thông: với các mô hình DL nhiều tham số, việc đồng bộ trọng số qua nhiều vòng học tiêu tốn băng thông và thời gian theo quy mô mô hình. Thứ hai là rủi ro suy diễn ngược dữ liệu từ cập nhật mô hình \cite{zhu2019dlg,zhao2020idlg,geiping2020inverting}. Thứ ba là mức độ phụ thuộc vào một mô hình toàn cục đồng nhất, khiến việc hỗ trợ các client dùng kiến trúc khác nhau trở nên khó khăn. Cũng chính từ ba giới hạn này mà nhóm nghiên cứu FD dần trở nên nổi bật hơn.

\subsection{Học máy liên kết dựa trên học chắt lọc}

Học máy liên kết dựa trên chắt lọc - Federated Distillation, thay thế dữ liệu được truyền trong quá trình học cộng tác. Thay vì gửi toàn bộ tham số mô hình, client gửi tri thức ở mức đầu ra hoặc biểu diễn trung gian, chẳng hạn logits, nhãn mềm (soft labels) hoặc prototype \cite{hinton2015kd,jeong2018federateddistillation,li2019fedmd,tan2022fedproto,li2024fdsurvey}. Sự thay đổi này đem lại một khác biệt lớn ở góc độ hệ thống: nếu học liên kết dựa trên trọng số yêu cầu các client phải có cùng kiến trúc mô hình DL, thì FD chỉ yêu cầu các client gặp nhau ở không gian dự đoán. Nhờ đó, các client có thể giữ kiến trúc riêng mà vẫn trao đổi tri thức với nhau một cách thống nhất.
\bigbreak\noindent
Trong nhóm này, Federated Distillation \cite{jeong2018federateddistillation} tổng hợp tri thức theo lớp, giúp lượng thông tin phải trao đổi phụ thuộc vào số lớp thay vì toàn bộ số tham số mô hình. FedMD \cite{li2019fedmd} khai thác một tập dữ liệu công khai để mọi client cùng dự đoán, từ đó xây dựng tín hiệu đồng thuận phục vụ distillation, và ưu điểm lớn nhất của FedMD là dễ mở rộng sang môi trường mô hình không đồng nhất nhờ tri thức được trao đổi ở dạng đầu ra trên cùng một tập mẫu. FedProto \cite{tan2022fedproto} lại đi theo hướng khác khi dùng prototype lớp, tức là biểu diễn đặc trưng trung bình, thay cho logit đầu ra. Cách tiếp cận này phù hợp khi hệ thống muốn chia sẻ thông tin ở mức biểu diễn nhưng vẫn giảm phụ thuộc vào tham số mô hình cụ thể.
\bigbreak\noindent
Trong lĩnh vực IDS, SSFL-IDS là một hướng nổi bật vì khai thác dữ liệu công khai không gán nhãn kết hợp với cơ chế bán giám sát để tạo pseudo label \cite{zhao2022ssfl}. Nhờ vậy, SSFL-IDS giảm được yêu cầu phải có sẵn một bộ dữ liệu công khai đã gán nhãn đầy đủ, vốn là một rào cản thực tế với nhiều hệ thống. Cùng hướng chọn lọc tri thức, FedSSD \cite{he2022fedssd} bổ sung cơ chế selective self-distillation để ưu tiên tri thức đáng tin hơn trong điều kiện Non-IID. Nhìn chung, các hướng này cho thấy FD không chỉ là giải pháp nén truyền thông, mà còn là cách thiết kế lại toàn bộ kênh hợp tác giữa các client.
\bigbreak\noindent
Về mặt ưu điểm cụ thể, cơ chế truyền thông dựa trên logit, hay logit-based communication, trước hết giúp chi phí truyền thấp hơn so với truyền mô hình, bởi một vector logit trên tập công khai có kích thước ổn định theo số lớp, trong khi tham số mô hình tăng theo chiều sâu và độ rộng của mạng. Khi số client lớn hoặc thiết bị nằm ở biên mạng, chênh lệch này trở thành một lợi thế vận hành cụ thể. Bên cạnh đó, không gian logit còn bảo toàn được mức độ tương đối giữa các lớp, giúp client nhận được nhiều thông tin hơn so với việc chỉ dựa vào nhãn cứng. Ngoài ra, một lợi ích không kém phần quan trọng là khả năng hỗ trợ mô hình không đồng nhất: một client có thể dùng dense, client khác dùng GRU, client khác nữa dùng kiến trúc sâu hơn, nhưng tất cả vẫn trao đổi tri thức được miễn là bài toán đầu ra được thống nhất.
\bigbreak\noindent
Tuy vậy, FD cũng không hoàn toàn hạn chế được nhiễu trong quá trình trao đổi tri thức. Nếu tập công khai không phù hợp, logit thu được sẽ kém giá trị; nếu public sample nằm quá xa phân phối riêng của nhiều client, tri thức tổng hợp từ những mẫu đó có thể gây nhiễu thay vì hỗ trợ. Tình huống này đặt ra hai yêu cầu bổ sung: cần cơ chế đánh giá public sample nào đáng dùng, và cần cơ chế đánh giá logit nào đáng tin trước khi tổng hợp. Đây là điểm nối trực tiếp đến các hướng dựa trên energy và EKD.

\subsection{Học máy liên kết phi tập trung}

Phần lớn các hệ thống FL và FD cổ điển đều giả định tồn tại một server trung tâm chịu trách nhiệm thu nhận thông tin từ client, tổng hợp rồi phân phối lại. Thiết kế đó giúp việc điều phối trở nên rõ ràng hơn, nhưng đồng thời cũng tạo ra nút cổ chai về niềm tin, hiệu năng và độ sẵn sàng. BrainTorrent và các khảo sát về FL phi tập trung (decentralized federated learning - DFL) đã nhấn mạnh rằng trong nhiều môi trường cộng tác thực tế, mô hình ngang hàng hoặc bán ngang hàng phản ánh điều kiện triển khai tốt hơn so với mô hình tập trung truyền thống \cite{braintorrent,beltran2023dflsurvey}.
\bigbreak\noindent
Trong IDS, giá trị của phi tập trung không chỉ nằm ở việc bỏ server. Quan trọng hơn, nó cho phép các bên tham gia duy trì quyền tự chủ cao hơn đối với dữ liệu, tiến trình học và hạ tầng nội bộ của mình. Với hệ thống liên tổ chức, điều này rất quan trọng vì không phải lúc nào cũng có một thực thể thứ ba được tất cả các bên tin cậy tuyệt đối. Các hướng như 2DF-IDS, các mô hình phi tập trung cho lightweight IDS và các biến thể anomaly detection ngang hàng đều cho thấy nghiên cứu về DFL-IDS đang mở rộng nhanh \cite{friha20232dfl-ids,nakip2023dfl-ids,attanayaka2023dfl-anomaly}. Tuy nhiên, đa số các hướng này vẫn chủ yếu trao đổi tham số mô hình, nên chưa giải quyết hết bài toán truyền thông và chưa khai thác trọn lợi thế của distillation trong môi trường mô hình không đồng nhất.
\bigbreak\noindent
Từ đó, có thể thấy một khoảng trống nghiên cứu cho một hệ thống vừa phi tập trung, vừa truyền tri thức nhẹ, vừa chấp nhận mô hình không đồng nhất, đồng thời vẫn thích nghi được với dữ liệu Non-IID và chất lượng tri thức không đồng đều giữa các client. Đây chính là mục tiêu nghiên cứu của khóa luận này.

\subsection{Học tiệm tiến trong môi trường liên kết}

Học tiệm tiến xử lý tình huống mô hình phải học dần các dữ liệu mới theo thời gian, và với IDS, giả định này gần với thực tế hơn nhiều so với giả định rằng toàn bộ lớp tấn công đều xuất hiện đầy đủ ngay từ đầu. Thay vào đó, dữ liệu mới liên tục xuất hiện và đưa chia theo từng tác vụ qua thời gian. Học tiệm tiến bao gồm ba kiểu chính: Domain-Incremental Learning, Task-Incremental Learning và Class-Incremental Learning. Domain-Incremental Learning là trường hợp các tác vụ có cùng không gian nhãn nhưng phân phối dữ liệu đầu vào thay đổi theo từng giai đoạn, ví dụ cùng phân loại “bình thường” và “tấn công” nhưng dữ liệu đến từ các môi trường mạng, thời điểm hoặc hệ thống khác nhau. Task-Incremental Learning là trường hợp dữ liệu được chia thành nhiều tác vụ riêng biệt, mỗi tác vụ có thể có tập nhãn khác nhau và mô hình thường được cung cấp thông tin về tác vụ khi dự đoán. Chẳng hạn, một tác vụ có thể là phân loại các dạng tấn công trong mạng doanh nghiệp, trong khi tác vụ khác là phân loại tấn công trong hệ thống IoT, và mô hình biết mẫu đầu vào thuộc tác vụ nào.

Trong khi đó, Class-Incremental Learning (CIL) là thiết lập thách thức và gần với thực tế hơn trong nhiều bài toán an ninh mạng. Ở đây, mô hình lần lượt học thêm các lớp mới theo thời gian, nhưng khi suy luận, mô hình phải phân biệt đồng thời giữa tất cả các lớp đã từng học mà không được biết mẫu dữ liệu thuộc giai đoạn hay tác vụ nào. Điều này phản ánh khá sát bài toán phát hiện và phân loại tấn công mạng: ban đầu hệ thống có thể chỉ nhận diện một số loại tấn công phổ biến như DoS, Probe hoặc Brute Force; sau đó, các kiểu tấn công mới như Botnet, Web Attack, Infiltration hoặc các biến thể zero-day dần xuất hiện. Khi đó, mô hình không chỉ cần học được các lớp tấn công mới mà còn phải tiếp tục nhận diện chính xác các lớp cũ và phân biệt chúng với lưu lượng bình thường. Vì vậy, Class-Incremental Learning không những là một vấn đề có độ khó cao mà còn là một bài toán thực tế quan trọng cho các hệ thống phát hiện xâm nhập hiện đại, nơi dữ liệu tấn công mạng liên tục mở rộng, phân bố thay đổi và yêu cầu mô hình duy trì năng lực nhận diện trên toàn bộ tập lớp đã học.

\bigbreak\noindent
Khó khăn chính của CIL là catastrophic forgetting, khi tối ưu theo dữ liệu mới làm tham số rời khỏi vùng từng biểu diễn tốt cho dữ liệu cũ. Để giải quyết catastrophic Forgetting, các phương pháp CIL thường dựa vào ba kỹ thuật chính: ràng buộc mô hình (regularization), học lại trên mẫu cũ đại diện (exemplar rehearsal) và kiến trúc động (dynamic architecture). Các kỹ thuật ràng buộc mô hình thường ràng buộc các trọng số quan trọng cho các tác vụ cũ như Elastic Weight Consolidation \cite{kirkpatrick2017ewc} và MAS (Memory Aware Synapses) \cite {aljundi2018mas}, hoặc ràng buộc đầu ra của mô hình thông qua học chắt lọc mô hình cũ hoặc biểu diễn đặc trưng (prototype) như Learning without Forgetting \cite{li2017lwf} và PASS \cite{zhu2021pass}. Mặt khác, một số phương pháp CIL thường lưu một số mẫu huấn luyện của tác vụ cũ và tái huấn luyện mô hình trong tác vụ mới, như Incremental classifier and representation learning (iCaRL) \cite{rebuffi2017icarl}, Bias Correction (BiC) \cite{wu2019bic}. Nhiều phương pháp CIL liên tục mở rộng kiến trúc mô hình phân loại để thích nghi với số nhãn tăng dần liên tục như FOSTER \cite{wang2022foster}.
\bigbreak\noindent
Khi CIL được đặt vào môi trường FL hoặc FD, bài toán không chỉ chịu Non-IID giữa client mà còn chịu thiếu lớp theo thời gian, nghĩa là mỗi client vừa học lệch về phân phối riêng, vừa không quan sát cùng một tiến trình xuất hiện lớp với các client khác. Trong thiết lập FCIL, mỗi client học theo một chuỗi tác vụ, còn tri thức giữa client chỉ có thể được hợp nhất thông qua trọng số, prototype, logits hoặc các thống kê biểu diễn, thay vì thông qua dữ liệu thô. Vì vậy, học liên kết tích hợp học tiệm tiến (Federated Class-Incremental Learning - FCIL) trở thành một bài toán kép: hệ thống phải đồng thời giải quyết catastrophic forgetting theo trục thời gian và lệch phân phối theo trục không gian. Đây là một vấn đề tương đối mới trong học máy liên kết, trong đó nghiên cứu đại diện là Glocal-Local Forgetting Compensation (GLFC) \cite{dong2022glfc}, và FEAT \cite{qi2026feat} là một nghiên cứu gần đây kết hợp replay và regularization cho bài toán FCIL.
\bigbreak\noindent
Đánh giá FCIL cũng cần nhìn rộng hơn một mô hình phân loại tĩnh. Bên cạnh accuracy, precision, recall và F1-score, hệ thống còn phải xét độ chính xác theo từng tác vụ, độ chính xác trung bình sau tác vụ cuối và forgetting measure giữa thời điểm học tốt nhất với thời điểm cuối. Nếu không có cơ chế cân bằng phù hợp, mô hình toàn cục rất dễ thiên lệch (bias) về những lớp hoặc những client xuất hiện gần đây, trong khi tri thức cũ dần suy yếu sau mỗi vòng hợp tác. Đây là lý do việc kết hợp FL, FD và CIL trong cùng một hệ thống IDS không chỉ mang ý nghĩa học thuật, mà còn đáp ứng nhu cầu thực tiễn.

\section{Mục tiêu, đối tượng và phạm vi nghiên cứu}

\subsection{Mục tiêu nghiên cứu}

Khóa luận hướng đến các mục tiêu sau:

\begin{itemize}
    \item Nghiên cứu cơ sở lý thuyết của Federated Learning, FD và CIL cho IDS.
    \item Xây dựng hệ thống mô phỏng Decentralized Federated Distillation kết hợp Class-Incremental Learning trên CICIoT2023.
    \item So sánh các thuật toán FL, FD và FCIL đã hiện thực trong mã nguồn, từ nhóm tổng hợp trọng số đến nhóm trao đổi tri thức mềm.
    \item Đánh giá ảnh hưởng của dữ liệu Non-IID, số lượng client, mô hình nền và thuật toán chắt lọc đến hiệu năng phát hiện.
    \item Đánh giá hệ thống đề xuất trong ngữ cảnh học tiệm tiến, với các phương pháp FCIL và phương CIL kết hợp FedAvg.
\end{itemize}

\subsection{Đối tượng nghiên cứu}

\begin{itemize}
    \item Dữ liệu phát hiện xâm nhập CICIoT2023.
    \item Mô hình DL cho phân loại lưu lượng mạng.
    \item Federated Learning và Federated Distillation.
    \item CIL trong môi trường tập trung và phân tán.
    \item Các thuật toán tổng hợp, chắt lọc và lưu giữ tri thức cũ.
\end{itemize}

\subsection{Phạm vi nghiên cứu}

Khóa luận tập trung vào mô phỏng FL và CIL trên dữ liệu tabular IDS. Nội dung không đi sâu vào triển khai một sản phẩm SIEM hoàn chỉnh hoặc một hệ thống vận hành thời gian thực. Các thí nghiệm được định hướng quanh tập dữ liệu CICIoT2023, các phân hoạch IID và Non-IID, cùng các thuật toán FL, chắt lọc tri thức và CIL đã được khảo sát.

\section{Phương pháp nghiên cứu}

Phương pháp nghiên cứu gồm bốn bước chính. Trước hết, khóa luận khảo sát lý thuyết và các công trình liên quan về IDS, FL, FD và CIL để xác định nền tảng của bài toán. Tiếp theo, mã nguồn hiện có được phân tích nhằm làm rõ kiến trúc pipeline, thuật toán huấn luyện, định dạng dữ liệu và cơ chế đánh giá. Từ đó, khóa luận này được thiết kế dưới dạng các thành phần: dữ liệu, client, mô hình cục bộ, cơ chế chắt lọc, vòng CIL và đánh giá. Cuối cùng, kế hoạch thực nghiệm được đề xuất nhằm so sánh các thuật toán theo khả năng phân loại của mô hình so với các phương pháp liên quan.

\section{Những điểm mới của đề tài}

\begin{itemize}
    \item Định hướng một hệ thống FD phi tập trung cho bài toán IDS, nhấn mạnh trao đổi tri thức thay vì trao đổi trọng số mô hình.
    \item Kết hợp CIL vào môi trường liên kết nhằm xử lý sự xuất hiện liên tục của lớp tấn công mới.
    \item Hệ thống hóa và so sánh nhiều nhóm thuật toán đã hiện thực, gồm FL, FD, học bán giám sát, học prototype và các phương pháp CIL.
    \item Tạo nền tảng thực nghiệm trên CICIoT2023 với dữ liệu Non-IID, nhiều client và nhiều kiến trúc mô hình.
\end{itemize}

\section{Cấu trúc Khóa luận tốt nghiệp}

Cấu trúc khóa luận gồm năm chương:

\begin{itemize}
    \item Chương~\ref{Chapter1}: \textbf{\nameref{Chapter1}} trình bày bối cảnh, lý do chọn đề tài, nghiên cứu liên quan và định hướng nghiên cứu.
    \item Chương~\ref{Chapter2}: \textbf{\nameref{Chapter2}} trình bày cơ sở lý thuyết về IDS, Federated Learning, chắt lọc tri thức, FD và CIL.
    \item Chương~\ref{Chapter3}: \textbf{\nameref{Chapter3}} trình bày phương pháp đề xuất và thiết kế hệ thống.
    \item Chương~\ref{Chapter4}: \textbf{\nameref{Chapter4}} trình bày hiện thực, thiết lập thực nghiệm, tiêu chí đánh giá và thảo luận kết quả.
    \item Chương~\ref{Chapter5}: \textbf{\nameref{Chapter5}} tổng kết khóa luận và đề xuất hướng phát triển.
\end{itemize}
